This note presents a measurement of dimuon angular correlations and yields 
  for muons arising from the decays of heavy-flavor hadrons in \PbPb\ 
  and $pp$ collisions at $\sqrtsnn=5.02$~\TeV.
The distributions of muon-pairs in azimuthal angle separation $\Delta\phi$, 
  are measured.
Combinatorial muon-pairs that arise from independently produced muons
  are estimated and subtracted from the measured distributions.
The correlated yields on the away-side ($\Delta\phi=\pi$) and widths of the 
  away-side correlation are measured as a function of the event centrality.
The modifications in the widths of the correlation are measured as 
  a function of centrality and compared to the widths measured in the \pp\ reference.
\iftoggle{IncludeYields}
{
 Similarly, the modifications in the pair-yields in Pb+Pb collisions, as a function of centrality,
  are compared to the \pp\ collisions. 
}
{}
Model calculations of such observables~\cite{Nahrgang:2013saa} indicate that 
  interactions of the heavy flavor quarks with the QGP produced in heavy-ion collisions,
  result in significant broadening of the away-side correlations in \Dphi\
  compared to the expectations from \pp\ collisions.
%These calculations also indicate that these measurements can help determine whether 
%  collisional or radiative processes dominate the energy loss of HF quarks as they 
%  traverse the QGP produced in heavy-ion collisions.
This implies that these measurements provide insight into the 
  question whether collisional or radiative processes dominate 
  the energy loss of HF quarks as they traverse the 
  QGP produced in heavy-ion collisions.

The \PbPb\ data used here were recorded by ATLAS in the 2015 and 2018 heavy ion runs
  and correspond to a total integrated luminosity of 1.94~\inb where
  0.50\inb\ was taken in 2015 and 1.44\inb\ in 2018.
The \pp\ data at \sqs=5.02~\TeV\ were recorded during a dedicated 
  reference run in 2017, and correspond to an integrated luminosity of 0.26~\ifb.
The analysis is performed on both same-charge muon pairs (i.e. both muons in the pair have the same sign)
  and opposite-sign muon pairs.
Most of the same-sign pairs are expected to arise from neutral $b$-osicllations, 
  and a smaller fraction from cascaded $b$-decays, where a $b$-hadron decays to 
  produce a $c$-hadron and a muon, and then the $c$-hadron decays to produce
  an opposite-sign muon.
So we expected similar behavior for the opposite charge pairs and same sign-pairs 
  when both come from $b$-$b$bar pairs. 
The $c$-$c$bar should contribute to only opposite-sign pairs (at leading order). 
So the same-sign and opposite sign distributions can help contrast 
  between muons-pairs from $b$-$b$bar and $c$-$c$bar.

The outline of the note is as follows:
Section~\ref{sec:datasets} describes the data, MC samples, triggers and 
  muon selections used in the analysis.
Section~\ref{sec:correction} presents the evaluation of the trigger 
  and reconstruction efficiencies, including some MC performance studies.
Section~\ref{sec:analysis} outlines the main steps of the analysis.
It describes the offline kinematic selections applied on the muons 
  and defines the observables for studying the dimuon correlations.
Section~\ref{sec:dpop_cuts} presents a study to estimate the fraction of
  background-pairs in the analysis, that arise from
  background muons produced from inflight decays of pions and Kaons
  and from secondary interactions in the calorimeter.
Section~\ref{sec:systematics} discusses the evaluation of the 
  systematic uncertainties for the measured observables.
Section~\ref{sec:results} shows the final results of the measurements, 
  summarizes the analysis and its main observations.

